% GENERATED, edit the French unit or the translation table
\mxnewpage\section*{EV4-A -- EAM preparation}\label{EV4-A}
\addcontentsline{toc}{section}{EV4-A -- EAM preparation}
\textbf{Paper reserved for the teacher}. Whole programme.\par Duration: 2 h. No calculator. Marked out of 20. Practice sessions in May and June; EAM date for the zone to be confirmed.
\subsection*{Part 1 -- Fluency (6 points)}
Twelve questions, \num{0,5} point each. Only one correct answer; no justification required. An incorrect or missing answer scores 0, with no loss of points.
\begin{enumerate}
\item Calculate $2+2\times3$.\par A. $11$\quad B. $7$\quad C. $8$\quad D. $12$
\item Calculate $\frac12+\frac13$.\par A. $\frac{5}{3}$\quad B. $\frac{1}{6}$\quad C. $\frac{5}{6}$\quad D. $\frac{2}{5}$
\item Calculate $\dfrac{2^2}{2^{0}}$.\par A. $2$\quad B. $1$\quad C. $7$\quad D. $4$
\item A value of 50 increases by \pc{20}. Final value?\par A. $70$\quad B. $40$\quad C. $60$\quad D. $10$
\item A value of 80 is obtained after a \pc{20} decrease. Initial value?\par A. $100$\quad B. $60$\quad C. $96$\quad D. $64$
\item A \pc{25} increase then a \pc{20} decrease: overall coefficient?\par A. $1$\quad B. $\frac{4}{5}$\quad C. $\frac{21}{20}$\quad D. $\frac{9}{10}$
\item The solution of $2x+2=0$ is\par A. $1$\quad B. $4$\quad C. $-1$\quad D. $-4$
\item The slope between $\pt{0;2}$ and $\pt{2;8}$ is\par A. $-3$\quad B. $3$\quad C. $6$\quad D. $2$
\item The mean of $2, 4, 6$ is\par A. $4$\quad B. $5$\quad C. $2$\quad D. $12$
\item A sector of $90^\circ$ represents what fraction of the total?\par A. $\frac{1}{4}$\quad B. $\frac{1}{3}$\quad C. $\frac{3}{4}$\quad D. $\frac{1}{2}$
\item A probability is $\dfrac{2}{5}$. Its complement is\par A. $\frac{5}{2}$\quad B. $\frac{3}{5}$\quad C. $\frac{2}{5}$\quad D. $\frac{1}{5}$
\item For $f(x)=x^2+2$, $f(-2)$ is\par A. $6$\quad B. $-2$\quad C. $4$\quad D. $8$
\end{enumerate}
\mxnewpage\subsection*{Part 2 -- Independent exercises (14 points)}
Every step must be justified. Exact results are expected unless stated otherwise.
\mxneed{12}
\subsection*{Exercise 1 -- Two linked variables (5 points)}
The random variable $X$ takes the values 1, 2 and 3, each with probability $\dfrac13$. On the same experiment, define the variable $Y=4-X$. We wish to check whether a rule valid for expectation remains true for variance.
\begin{enumerate}
\item Calculate $E(X)$ and $E(Y)$. \hfill(1 pt)
\item Calculate $V(X)$, $V(Y)$ and $V(X+Y)$. \hfill(2 pts)
\item Use these results to refute the statement: ``The variance of a sum is always equal to the sum of the variances.'' \hfill(2 pts)
\end{enumerate}
\mxneed{12}
\subsection*{Exercise 2 -- Comparing a curve and a line (5 points)}
The functions $f(x)=e^{-x}$ and $g(x)=1-x$ are defined on $\R$. Let $h=f-g$. Comparing their curves will then allow us to check the direction of an approximation.
\begin{enumerate}
\item Differentiate their difference. \hfill(1 pt)
\item Prove $f(x)\geq g(x)$ for every real number x. \hfill(2 pts)
\item Write the affine approximation of $e^{-\num{0,02}}$ and state whether it is from below or from above. \hfill(2 pts)
\end{enumerate}
\mxneed{12}
\subsection*{Exercise 3 -- Normal vector and radius (4 points)}
In an orthonormal frame, consider the point $I\pt{1;1}$ and the line $d$ with equation $3x+4y-17=0$. Let $H$ be the orthogonal projection of $I$ onto $d$. We look for the circle with centre $I$ tangent to $d$.
\begin{enumerate}
\item Calculate the distance from I to d. \hfill(1 pt)
\item Calculate the projection H. \hfill(1 pt)
\item Give the equation of the circle with centre I tangent to d and check that H belongs to the circle. \hfill(2 pts)
\end{enumerate}
